\documentclass[12pt,a4paper]{article}

\usepackage[utf8]{inputenc}
\usepackage[T1]{fontenc}
\usepackage[brazil]{babel}
\usepackage[margin=2.5cm]{geometry}
\usepackage{amsmath}
\usepackage{amssymb}
\usepackage{array,booktabs,needspace,tabularx}
\usepackage{float,graphicx}
\graphicspath{{./}{dba-studies/projeto-pesquisa/}}
\usepackage[authoryear,round]{natbib}
\usepackage{hyperref}
\hypersetup{hidelinks}
\let\cite\citep

\setlength{\parskip}{0.5em}
\setlength{\parindent}{0pt}
\widowpenalty=10000
\clubpenalty=10000

\title{AdaptGov: aprendizado de política fiscal adaptativa em um modelo macroeconômico baseado em agentes}
\author{Daniel Andrade}
\date{Outrubro de 2026}

\begin{document}

\maketitle
\vspace{-3em}

\section{Introdução, objetivos e justificativa}

\subsection{Introdução}

Um sistema complexo se caracteriza por possuir muitos agentes interagindo de forma a produzir comportamentos emergentes, ou seja, que não podem ser trivialmente derivados a partir de seus comportamentos individuais \cite{newman2011}. Estes fenômenos emergentes por sua vez afetam os agentes individualmente, criando retroalimentações ("feedbacks") positivas e negativas.

A economia é um sistema complexo onde bilhões de agentes heterogêneos, entre trabalhadores, empresas, bancos e instituições governamentais, interagem em compras e vendas, empréstimos, pagamentos, investimentos, resultando em fenômenos emergentes como preços, inflação e crescimento. Além disso, esses atores aprendem e evoluem ao longo do tempo, o que modifica fundamentalmente a estrutura do sistema, tornando-o adaptativo \cite{arthur2013}. 

A aplicação da lente teórica de sistemas complexos adaptativos para o estudo da economia tem se mostrado uma alternativa promissora ao paradigma predominante da teoria econômica, o neoclássico, que se ancora em pressupostos teóricos pouco críveis, como a existência de um equilíbrio geral, as expectativas racionais e os agentes representativos \cite{gallegati2012}.

Nesse sentido, o campo da economia complexa (“complexity economics”), seria uma formulação mais geral da economia neoclássica, pois estuda as estruturas econômicas a partir de premissas menos restritas \cite{arthur2013}. 

Os modelos baseados em agentes ("ABMs") tem se constituído uma ferramenta poderosa para o estudo da economia complexa através de ferramentas computacionais \cite{dawid2018}. Neles, os agentes são simulados por meio de regras comportamentais que regem tanto seu comportamento individual quanto suas interações, sem a necessidade de um controle central. As variáveis agregadas resultantes emergem como resultado computacional dos comportamentos e interações ao longo do tempo, sem a utilização de pressupostos de equilíbrio e de informação e racionalidade perfeita.

O trabalho proposto por Poledna et. al \cite{poledna2023}, é um modelo baseado em agentes para a macroeconomia, ou seja, focado no estudo de propriedades agregadas da economia, como inflação, PIB, crescimento, desemprego \cite{dawid2018}, em particular para a Áustria, que é um país pequeno, industrial e com economia aberta. Ele é ancorado nas estatísticas nacionais e em regras autoregressivas de primeira ordem, realizando previsões macroeconômicas comparávais aos modelos DSGE ("Dynamic Stochastic General Equilibrium"), que representa do paradigma atual da macroeconomia neoclássica. O BeforeIT.jl \cite{glielmo2025} oferece uma implementação aberta e extensível deste modelo na linguagem Julia feita pela equipe de tecnologia da informação do Banco da Itália. Nos trabalhos apresentados, a política fiscal é constituída majoritariamente como parâmetros fixos que delimitam impostos, benefícios sociais e gastos do governos. Nesse sentido, há espaço para o estudo de uma política fiscal variável e adaptativa, ou seja, que reage ao estado do modelo. 

No campo dos ABMs para macroeconomia, novos trabalhos tem alavancado em inteligência artificial ("IA"), como em \cite{zheng2022} e \cite{brusatin2024}, onde técnicas de aprendizado por reforço são utlizadas para treinar agentes com comportamentos racionais e avaliar o resultado macroeconômico. Em \cite{li2024} LLMs ("Large Language Models") fazem o papel de consumidores e trabalhadores em um modelo que busca reproduzir qualitativamente fatos estilizados macroeconômicos.  

Por outro lado, a análise de sensibilidade e os metas modelos tem se constituído ferramentas importante para a avaliação e emulação de ABMs em diversas áreas, pois podem auxiliar na robustez e explicabilidade aos resultados dos modelos em questão como pode ser visto em \cite{angione2022}. 

Por ser um modelo aberto, extensível, calibrado para as estatísticas nacionais e com desempenho preditivo similar aos DSGEs nos testes dos autores, o Before.IT pode ser um campo fértil para investigar políticas fiscais em um modelo macroeconômico baseado em agentes. Este trabalho propõe realizar análises de sensibilidade dos parâmetros fiscais e treinar com aprendizagem por reforço uma política fiscal adaptativa. Essa política será comparada a configurações fixas de parâmetros com relação a indicadores de bem-estar (e.g. considerando PIB, desemprego e inflação), dívida pública e estabilidade macroeconômica. Assim, busca-se avaliar se a adaptação da política fiscal melhora o resultado das simulações e em quais casos. Esse tipo de exploração pode ser um passo importante para tornar o BeforeIT.jl e outros ABMs macroeconômicos úteis ao auxiliar tomadores de decisões governamentais na formulação de políticas públicas. 

\subsection{Perguntas de Pesquisa}

\textbf{Pergunta central}

No modelo BeforeIT.jl calibrado para a Áustria, uma política fiscal adaptativa treinada com aprendizado por reforço proporciona maior bem-estar que políticas de parametros fixos, considerando referências fiscais de dívida e défict inspiradas nas regras europeias?

\textbf{Perguntas secundárias}
\begin{itemize}
    \item Quais alterações nos parâmetros da política fiscal produzem resultados estatisticamente significantes?
    \item Em quais cenários os resultados de uma política fiscal adaptativa superam as regras fixas e políticas fiscais?
\end{itemize}

\subsection{Justificativa}

Modelos baseados em agentes são uma ferramenta poderosa para a análise da macroeconomia com o potencial de informar a sociedade e agencias governamentais sobre o comportamento passado e futuro de indicadores chaves como inflação, desemprego, crescimento \cite{arthur2013,gallegati2012,dawid2018}. 

Nesse cenário, a construção de análise de sensibilidade é uma ferramenta importantes para avaliar a validade, explicabilidade e robustez dos modelos em questão. 

Além disso, os modelos baseados em agentes tem tido sua fronteira explorada através da utilização de inteligência artificial para a construção de agentes adaptativos que aprendem ao longo da simulação\cite{angione2022, zheng2022}. 

Este trabalho busca contribuir para o campo dos modelos macroeconomico baseados em agentes ao explorar um modelo de fronteira, construindo, validando e interpretando políticas fiscais adaptativa construídas através de aprendizagem por reforço. Assim, busca-se extender a aplicação do modelo de um papel preditivo para também um papel informativo sob eventuais políticas fiscais a serem aplicadas por um formulador de política economica.  

\subsection{Objetivos}

\Needspace{5\baselineskip}
Objetivo geral:

Desenvolver e avaliar políticas fiscais adaptativas no modelo Before.IT, calibrado para Áustria, através de técnicas de análise de sensibilidade e aprendizagem por reforço, comparando-as com políticas de parametros fixos, quanto a bem-estar, dívida pública e estabilidade macroeconômica.

Objetivos específicos:
\begin{itemize}
    \item Formular o problema de política fiscal no contexto do modelo Before.IT, com seus parâmetros, variáveis livres e restrições. 
    \item Realizar análise de sensibilidade nas variáveis de estudo do problema de fiscal quantificando seus efeitos nos indicadores de avaliação selecionados.
    \item Propor, treinar e avaliar um modelo de aprendizado por reforço para a política fiscal.
    \item Avaliar e interpretar os resultados considerando bem-estar, dívida pública e estabilidade macroecoznômica. 
\end{itemize}

\section{Revisão bibliográfica}

Uma revisão geral sobre o campo de estudo dos sistemas complexos pode ser econtrada em \cite{newman2011}, que os conceitua como sistemas onde muitos atores interagem formando comportamentos que não podem ser explicados por suas ações individuais, chamados de emergentes. 

Em \cite{arthur2013}, \cite{gallegati2012} e \cite{dawid2018} é possível encontrar uma conexão entre sistemas complexos e as ciências economicas. Em especial, a abordagem de sistemas complexos para a economia é colocada como visão não reducionista e não axiomática, onde se estuda a interações economicas deses agentes. 

A modelagem baseada em agentes para economia ("Agent-based Computational Economics") é definida em \cite{dawid2018} como o estudo computacional de processos economicos modelados a partir de sistemas dinamicos de agentes interativos. Além disso, \citet{dawid2018} descrevem um estudo de \citet{dosi2015} com regras fiscais inspiradas no Pacto de Estabilidade e Crescimento e no Pacto Fiscal Europeu, que usam como referências um déficit de 3\% do PIB e uma dívida de 60\% do PIB. Neste estudo, esses valores servirão de referência para a função de recompensa e para a avaliação da política fiscal adaptativa.

Em \cite{poledna2023} foi construído um ABM de previsão macroeconômica para a economia da Áustria. Em \cite{glielmo2025} esse trabalho foi adaptado para a linguagem julia e disponibilidado de maneira open source. 

Em \cite{zheng2022}, \cite{zheng2022} e \cite{brusatin2024} é possível encontrar exemplos da aplicação de ferramentas de IA, como LLMs e aprendizagem por reforço, para aplicações em ABMs macroeconômicos.

\section{Modelo de referência}

As principais motivações do modelo BeforeIT.jl são:

\begin{itemize}
    \item Realizar `previsões macroeconômicas` utilizando modelagem baseada em agentes.
    \item Construir uma modelagem verossimilhante, calibrando o modelo de acordo com as estatísticas governamentais e executando a simulação com escala 1:1 da população do país.
\end{itemize}

O modelo inicial foi calibrado para a economia da Áustria, que é um país pequeno e altamente integrado com a economia externa ($ \approx 9M$ habitantes  e $\approx 50\%  $ do PIB em importações/exportações). Os agentes do modelo são famílias ativas e invativas, empresas e empresários, governo, um banco comercial, um banco central e empresas extrangeiras.

\Needspace{10\baselineskip}
Cada etapa temporal do modelo equivale a um `trimestre`. As etapas-chave que levam o modelo de $t$ para $t+1$ são:

\begin{enumerate}
    \item \textbf{Formação de expectativas coletivas:} nesta etapa é realizada uma regressão de primeira ordem para a previsão do PIB e da inflação. Essas previsões são utilizadas pelos diversos agentes, como empresas para estimar a produção do próximo ciclo e consumidores para estimar a renda do próximo ciclo. 
    \item \textbf{Execução dos choques exógenos:} nesta etapa são sorteados choques exógenos que modelam demanda e oferta do resto do mundo, ou seja, que quantificam as importações e exportações, bem como a inflação da zona do euro. 
    \item \textbf{Execução dos mercados de crédito e trabalho:} nesta etapa as empresas, após determinarem seus planos de produção, determinam o quanto querem de crédito do banco comercial. O banco então fornece crédito em uma ordem aleatória às empresas, obedecendo as regras regulatórias de fornecimento de crédito. Depois, as empresas determinam o qual será a força de trabalho a ser contrada ou demitida, que é selecionada aleatóriamente. 
    \item \textbf{Execução dos mercados de bens:} inicialmente o mercado intra-empresas é executado através do mencanismo de busca e pareamento, onde, para cada bem desejado, a empresas consumidora é selecionada aleatoriamente e os fornecedores são selecionados de com probabilidade proporcional ao seu preço e tamanho. O mesmo processo é então repetido para os consumidores finais e empresas de bem de consumo. 
    \item \textbf{Atualização da contabilidade:} são atualizados os preços, rendas, saldos, impostos e agregados econômicos. 
\end{enumerate}

\begin{figure}[H]
\centering
\includegraphics[width=\linewidth]{novo-fluxo-completo-mermaid-diagram.png}
\caption{Fluxo de um período de simulação do modelo de referência.}
\label{fig:novo-fluxo-completo-mermaid-diagram}
\end{figure}

\clearpage

\section{Metodologia e plano de atividades}

A pesquisa será baseada em simulações computacionais do modelo BeforeIT.jl calibrado para a Áustria. Os resultados da simulação no tempo \(t\), representados por \(\mathbf{y}_{t}\), dependem dos parâmetros \(\theta\), do vetor de variáveis de estado
iniciais \(\mathbf{x}_0\) e dos choques aleatórios \(\varepsilon_{t}\):

\[
\mathbf{x}_{t+1}=F(\mathbf{x}_t,\theta,\varepsilon_{t+1}),
\qquad
\mathbf{y}_t=g(\mathbf{x}_t).
\]

O agente governamental adaptativo terá como objetivo maximizar a soma da sua função de recompensa descontada ao longo do tempo: 

\[
S = \sum_{t=1}^{t_{\mathrm{sim}}} \gamma^t r_{t}
\]

Assim, a recompensa em cada etapa temporal \(r_{t}\) será uma funcão do espaço de estados do modelo \(y_{t}\). A modelagem de aprendizado por reforço será baseada na proposta de \cite{angione2022}. Será utilizado o algoritmo padrão de Q-learning, com uma política \(\varepsilon\) greedy. O agente terá um espaço de estado \(\mathcal{S}\) e espaço de ação discretos \(\mathcal{A}\), e aprenderá a função de valor \(Q : \mathcal{S} \times \mathcal{A} \to \mathbb{R}\), utilizando a equação de Bellman para estimar sua função de recompensa \(r\) dado um espaço de estado e possíveis ações \((s,a)\). 

O agente governamental terá como seu espaço de estados observações dos prinicipais agregados econômicos de interesse para a política fiscal, como função do espaço de estados geral do modelo \(\mathbf{y}_{t}\). A saber: 

\[
\mathbf{s}_t = f(\mathbf{y}_t)
= \left(Y_t^{\mathrm{real}},\, u_t,\, \pi_t,\, d_t,\, b_t\right).
\]


\begin{itemize}
    \item \(Y_t^{\mathrm{real}}\): PIB real no trimestre \(t\);
    \item \(u_t\): taxa de desemprego;
    \item \(\pi_t\): taxa de inflação;
    \item \(d_t\): déficit público acumulado nos últimos quatro trimestres
    dividido pelo PIB nominal acumulado no mesmo período;
    \item \(b_t\): dívida pública ao fim do trimestre \(t\) dividida pelo
    PIB nominal anualizado a partir de até quatro trimestres;
\end{itemize}

Esse espaço de estado será discretizado a para o treinamento Q-learning, utilizando \(n_S\) divisões espaçadas entre um mínimo \(S_{\min}\) e um máximo \(S_{\max}\). Cada valor real de observação será aproximado ao valor discreto mais próximo. 

Para o espaço de ações, foram selecionados o cojunto dos parâmetros \(\theta\) que interessam a política fiscal e que tipicamente estão sob controle da política fiscal, ou seja, dos impostos e gastos públicos. A saber:

\[
p_t =
\left(
\tau_t^{\mathrm{INC}},
\tau_t^{\mathrm{VAT}},
\theta_t^{\mathrm{UB}},
\log m_t^G
\right)
\]
\begin{itemize}
    \item \(p_t\): vetor de parâmetros da política fiscal no trimestre \(t\);
    \item \(\tau_t^{\mathrm{INC}}\): alíquota do imposto de renda no trimestre \(t\);
    \item \(\tau_t^{\mathrm{VAT}}\): alíquota do imposto sobre valor agregado (IVA) no trimestre \(t\);
    \item \(\theta_t^{\mathrm{UB}}\): taxa de reposição do benefício de desemprego no trimestre \(t\);
    \item \(\log m_t^G\): logaritmo do multiplicador da demanda de consumo do governo no trimestre \(t\);
    \item \(\mathcal{P}\): conjunto de combinações de política fiscal admissíveis.
\end{itemize}


Após observar o estado \(s_t\), o agente governamental escolhe
um vetor de ajustes
\[
a_t =
\left(
a_t^{\mathrm{INC}},
a_t^{\mathrm{VAT}},
a_t^{\mathrm{UB}},
a_t^G
\right) \in \mathcal{A}.
\]
Esses ajustes determinam os parâmetros da política fiscal do
trimestre seguinte:
\[
p_{t+1}
=
\left(p_t+a_t\right),
\qquad
p_t =
\left(
\tau_t^{\mathrm{INC}},
\tau_t^{\mathrm{VAT}},
\theta_t^{\mathrm{UB}},
\log m_t^G
\right).
\]

Os demais parâmetros do modelo serão mantidos os mesmos da calibração do BeforeIT.jl.

Em seguida será treinada a política fiscal adaptativa utilizando apredizagem por reforço. Essa política terá como objetivo maximixar o bem-estar, minimizando o deficit público e a dídiva pública de acordo com a política europeia de estabilização e crescimento, através da seguinte função de recompensa:
\[
\begin{aligned}
r_t={}&\alpha\log\left(\frac{Y_t^{\mathrm{real}}}{Y_0^{\mathrm{real}}}\right)
-\beta u_t-\gamma|\pi_t-\pi^*|\\
&-\lambda_a\left\|p_t-p_{t-1}\right\|_2^2
-\lambda_d\left[\max\left(0,d_t-0{,}03\right)\right]^2\\
&-\lambda_b\left[\max\left(0,b_t-b_t^{\mathrm{meta}}\right)\right]^2.
\end{aligned}
\]

\begin{itemize}
    \item \(r_t\): recompensa no trimestre \(t\);
    \item \(Y_t^{\mathrm{real}}\): PIB real no trimestre \(t\);
    \item \(Y_0^{\mathrm{real}}\): PIB real inicial;
    \item \(u_t\): taxa de desemprego;
    \item \(\pi_t\): taxa de inflação;
    \item \(\pi^*\): meta de inflação;
    \item \(p_t\): vetor de parâmetros da política fiscal aplicado
    no trimestre \(t\);
    \item \(p_{t-1}\): vetor aplicado no trimestre anterior;
    \item \(d_t\): déficit público acumulado nos últimos quatro trimestres
    dividido pelo PIB nominal acumulado no mesmo período; nos primeiros
    trimestres, usa-se a janela disponível;
    \item \(b_t\): dívida pública ao fim do trimestre \(t\) dividida pelo
    PIB nominal anualizado a partir de até quatro trimestres;
    \item \(b_0\): razão dívida/PIB no início da simulação;
    \item \(b_t^{\mathrm{meta}}=0{,}60+\max(0,b_0-0{,}60)
    \left(\frac{19}{20}\right)^{t/4}\): trajetória que reduz em 1/20 por
    ano a distância inicial da dívida à meta de 60\% do PIB;
    \item \(\alpha,\beta,\gamma,\lambda_a>0\): pesos dos componentes
    da recompensa.
    \item \(\lambda_d,\lambda_b>0\): pesos das penalidades por excesso de
    déficit e desvio da trajetória da dívida, respectivamente.
\end{itemize}



Inicialmente, será realizada uma análise de sensibilidade dos
parâmetros da política fiscal através de simulações de Monte Carlo. Serão selecionadas \(N\) configurações por amostragem por hipercubo latino ("LHS").
Cada configuração será mantida fixa durante 12 trimestres, e seus
resultados serão calculados com \(J\) sementes aletórias comuns. Os pares
formados pelos parâmetros e pelos resultados médios serão usados
para estimar a influência de cada alavanca fiscal.

\[
\{p_i\}_{i=1}^{N}=\mathrm{LHS}(\mathcal{P},N),
\qquad
\bar{\mathbf{y}}_{i,12}
=\frac{1}{J}\sum_{j=1}^{J}
S_{12}(p_i,\mathbf{x}_0,\varepsilon_j).
\]

\begin{itemize}
    \item \(p_i\): vetor de parâmetros da configuração \(i\);
    \item \(N\): número de configurações amostradas;
    \item \(J\): número de sementes por configuração;
    \item \(\mathbf{x}_0\): estado inicial do modelo;
    \item \(\varepsilon_j\): choques e sorteios associados à semente \(j\);
    \item \(S_{12}\): resultado do simulador após 12 trimestres;
    \item \(\bar{\mathbf{y}}_{i,12}\): vetor médio de resultados da
    configuração \(i\).
\end{itemize}

A análise de sensibilidade será uma primeira etapa informativa sobre o potencial impacto da política fiscal nas variáveis de interesse do modelo como inflação, crescimento do PIB, desemprego, deficit público e dívida do governo. Caso ela mostre que alterar os parâmetros da política fiscal não altera significativamente as variáveis de interesse, a política adaptativa pode se mostrar inóqua.  




Os pesos da função de recompensa serão definidos a partir da análise de sensibilidade, antes do treinamento da política adaptativa. 

\[
\begin{aligned}
R_{i,j,T}&=\sum_{t=1}^{T}r_t^{(i,j)},\\
D_{i,j,T}&=\sum_{t=1}^{T}
\left[\max\left(0,d_t^{(i,j)}-0{,}03\right)\right]^2,\\
B_{i,j,T}&=\sum_{t=1}^{T}
\left[\max\left(0,b_t^{(i,j)}-b_t^{\mathrm{meta},(i,j)}\right)\right]^2,\\
\bar{\mathbf{z}}_{i,T}&=\frac{1}{J}\sum_{j=1}^{J}
\begin{pmatrix}R_{i,j,T}\\D_{i,j,T}\\B_{i,j,T}\end{pmatrix}.
\end{aligned}
\]

\[
\hat{\mathbf{z}}_T(p)=\hat h(p),\qquad
\hat h=\mathop{\mathrm{arg\,min}}\limits_{h\in\mathcal{H}}
\sum_{i=1}^{N}\left\|\bar{\mathbf{z}}_{i,T}-h(p_i)\right\|_2^2.
\]

\[
p^\star=\mathop{\mathrm{arg\,max}}\limits_{p\in\mathcal{P}}
\left[\hat{\mathbf{z}}_T(p)\right]_1,\qquad
p_0=p^\star,\qquad
p_t=\mu(\mathbf{x}_{t-1})=p^\star+\Delta_\phi(\mathbf{x}_{t-1})\in\mathcal{P}.
\]

Todos os procedimentos destacados serão primeiramente realizados em uma escala reduzida do modelo (1:1000) e posteriormente serão reproduzidos e validados na escala 1:1 para a economia da Áustria.


\begin{table}[H]
\centering
\caption{Plano de atividades por semestre}
\label{tab:plano-atividades}
\footnotesize
\setlength{\tabcolsep}{3pt}
\renewcommand{\arraystretch}{1.15}
\begin{tabularx}{\textwidth}{@{}>{\raggedright\arraybackslash}X >{\raggedright\arraybackslash}p{0.22\textwidth} cccc@{}}
\toprule
\textbf{Atividade e método} & \textbf{Entregável} & \textbf{Sem 1} & \textbf{Sem 2} & \textbf{Sem 3} & \textbf{Sem 4} \\
\midrule
Revisão de literatura e documentação de funcionamento do modelo
& Lacunas de pesquisa definidas & X & & & \\
\addlinespace[2pt]
Realizar análise de sensibilidade nos parâmetros fiscais do modelo
& Intervalos viáveis de ajuste na política fisscal e seus respectivos efeitos & X & X & & \\
\addlinespace[2pt]
Definir objetivo da política fiscal
& Critérios de avaliação definidos & & X & & \\
\addlinespace[2pt]
Avaliar o paper do aprendizado de máquina supervisionado e desenvolver agente governamental usando aprendizagem por reforço 
& Treinar política e documentar decisões & & & X & \\
\addlinespace[2pt]
Testar políticas em novas simulações, com e sem choques economicos, comparando resultados e incertezas
& Resultados de bem-estar, fiscais e robustez & & & X & X \\
\addlinespace[2pt]
Consolidar resultados, qualifcar e finalizar escrita 
& Qualificação e defesa & X & X & & X \\
\bottomrule
\end{tabularx}
\end{table}

\clearpage
\section{Resultados parciais, se houver}

\clearpage
\renewcommand{\refname}{Referências bibliográficas}
\begin{thebibliography}{99}

\bibitem[Angione et al.(2022)]{angione2022} ANGIONE, Claudio; SILVERMAN, Eric; YANESKE, Elisabeth. Using machine learning as a surrogate model for agent-based simulations. \emph{PLOS ONE}, v. 17, n. 2, e0263150, 2022. DOI: \url{https://doi.org/10.1371/journal.pone.0263150}.

\bibitem[Arthur(2013)]{arthur2013} ARTHUR, W. B. \emph{Complexity Economics: A Different Framework for Economic Thought}. Santa Fe Institute Working Paper, 2013. Disponível em: \url{https://sites.santafe.edu/~wbarthur/Papers/Comp.Econ.SFI.pdf}.

\bibitem[Brusatin et al.(2024)]{brusatin2024} BRUSATIN, Simone; PADOAN, Tommaso; COLETTA, Andrea; DELLI GATTI, Domenico; GLIELMO, Aldo. Simulating the Economic Impact of Rationality through Reinforcement Learning and Agent-Based Modelling. In: \emph{Proceedings of the 5th ACM International Conference on AI in Finance}. ACM, 2024. p. 159--167. DOI: \url{https://doi.org/10.1145/3677052.3698621}.

\bibitem[Dawid e Delli Gatti(2018)]{dawid2018} DAWID, Herbert; DELLI GATTI, Domenico. Agent-Based Macroeconomics. In: \emph{Handbook of Computational Economics}. Elsevier, 2018. v. 4, p. 63--156. DOI: \url{https://doi.org/10.1016/bs.hescom.2018.02.006}.

\bibitem[Dosi et al.(2015)]{dosi2015} DOSI, Giovanni; FAGIOLO, Giorgio; NAPOLETANO, Mauro; ROVENTINI, Andrea; TREIBICH, Tania. Fiscal and monetary policies in complex evolving economies. \emph{Journal of Economic Dynamics and Control}, v. 52, p. 166--189, 2015. DOI: \url{https://doi.org/10.1016/j.jedc.2014.11.014}.

\bibitem[Gallegati e Kirman(2012)]{gallegati2012} GALLEGATI, Mauro; KIRMAN, Alan. Reconstructing Economics: Agent Based Models and Complexity. \emph{Complexity Economics}, v. 1, n. 1, p. 5--31, 2012. DOI: \url{https://doi.org/10.7564/12-COEC2}.

\bibitem[Glielmo et al.(2025)]{glielmo2025} GLIELMO, Aldo; DEVETAK, Mitja; MELIGRANA, Adriano; POLEDNA, Sebastian. BeforeIT.jl: High-Performance Agent-Based Macroeconomics Made Easy. \emph{arXiv}, arXiv:2502.13267, 2025. DOI: \url{https://doi.org/10.48550/arXiv.2502.13267}. Preprint.

\bibitem[Li et al.(2024)]{li2024} LI, Nian; GAO, Chen; LI, Mingyu; LI, Yong; LIAO, Qingmin. EconAgent: Large Language Model-Empowered Agents for Simulating Macroeconomic Activities. \emph{arXiv}, arXiv:2310.10436, 2024. DOI: \url{https://doi.org/10.48550/arXiv.2310.10436}. Preprint.

\bibitem[Newman(2011)]{newman2011} NEWMAN, M. E. J. \emph{Complex Systems: A Survey}. arXiv:1112.1440, 2011. Disponível em: \url{https://arxiv.org/abs/1112.1440}.

\bibitem[Poledna et al.(2023)]{poledna2023} POLEDNA, Sebastian; MIESS, Michael Gregor; HOMMES, Cars; RABITSCH, Katrin. Economic forecasting with an agent-based model. \emph{European Economic Review}, v. 151, art. 104306, 2023. DOI: \url{https://doi.org/10.1016/j.euroecorev.2022.104306}.

\bibitem[Zheng et al.(2022)]{zheng2022} ZHENG, Stephan; TROTT, Alexander; SRINIVASA, Sunil; PARKES, David C.; SOCHER, Richard. The AI Economist: Taxation policy design via two-level deep multiagent reinforcement learning. \emph{Science Advances}, v. 8, n. 18, eabk2607, 2022. DOI: \url{https://doi.org/10.1126/sciadv.abk2607}.

\end{thebibliography}

\end{document}
